\documentclass[11pt]{article}

\usepackage[margin=1in]{geometry}
\usepackage{amsmath, amssymb}
\usepackage{booktabs}
\usepackage{siunitx}
\usepackage{hyperref}
\usepackage{algorithm}
\usepackage{algpseudocode}

\sisetup{per-mode=symbol}

\newcommand{\vect}[1]{\mathbf{#1}}
\newcommand{\dd}{\,\mathrm{d}}
\newcommand{\Rp}{R_p}
\newcommand{\Ra}{R_a}

\title{Single-Scattering Atmospheric Rendering\\
       \large Formulas for a planetary atmosphere in a ray tracer}
\author{Connor Petri\\CS-116A}
\date{\today}

\begin{document}
\maketitle

\section{Overview}

The atmosphere is modelled as a spherical shell of gas between the planet
surface at radius $\Rp$ and the top of the atmosphere at radius $\Ra$, both
centred on the planet centre $\vect{c}$. Sunlight entering the shell is
scattered by air molecules (Rayleigh scattering) and by aerosols (Mie
scattering). The \emph{single-scattering} approximation assumes each photon
reaching the camera has scattered exactly once; light that scatters two or
more times is ignored.

For every camera ray the renderer computes two quantities:
\begin{enumerate}
  \item the \textbf{in-scattered radiance} $L_{\mathrm{in}}$, sunlight redirected
        toward the camera by the air along the ray, and
  \item the \textbf{transmittance} $T$, the fraction of light from whatever the
        ray hits (the planet surface or space) that survives the trip to the
        camera.
\end{enumerate}
The final pixel colour is
\begin{equation}
  L = T \, L_{\mathrm{surface}} + L_{\mathrm{in}} .
  \label{eq:composite}
\end{equation}

\section{Density profile}

Let $h(\vect{x}) = \lVert \vect{x} - \vect{c} \rVert - \Rp$ be the altitude of a
point above the surface. Both particle types have a density that falls
exponentially with altitude, normalised to $1$ at sea level:
\begin{equation}
  \rho_R(h) = \exp\!\left(-\frac{h}{H_R}\right), \qquad
  \rho_M(h) = \exp\!\left(-\frac{h}{H_M}\right),
  \label{eq:density}
\end{equation}
where $H_R$ and $H_M$ are the Rayleigh and Mie \emph{scale heights}.

\section{Scattering coefficients}

The scattering coefficient $\beta$ is the fraction of light scattered per unit
distance at sea level. For Rayleigh scattering it depends on wavelength
$\lambda$ as
\begin{equation}
  \beta_R(\lambda) = \frac{8\pi^3 \left(n^2 - 1\right)^2}{3 N \lambda^4},
  \label{eq:beta_rayleigh}
\end{equation}
with $n$ the refractive index of air and $N$ the molecular number density at
sea level. The $\lambda^{-4}$ dependence is why blue light (short $\lambda$)
scatters roughly five times more strongly than red. In practice $\beta_R$ is
treated as a per-channel RGB constant. Mie scattering is approximately
wavelength independent, so $\beta_M$ is a single scalar.

\section{Optical depth and transmittance}

The \emph{optical depth} of a path $\gamma$ is the density integrated along it:
\begin{equation}
  \tau_R(\gamma) = \int_\gamma \rho_R\big(h(\vect{x})\big) \dd s, \qquad
  \tau_M(\gamma) = \int_\gamma \rho_M\big(h(\vect{x})\big) \dd s .
  \label{eq:optical_depth}
\end{equation}
By the Beer--Lambert law, the fraction of light that traverses $\gamma$ without
being scattered or absorbed is
\begin{equation}
  T(\gamma) = \exp\!\Big( -\big[\beta_R \, \tau_R(\gamma) + \beta_M^{e} \, \tau_M(\gamma)\big] \Big),
  \label{eq:transmittance}
\end{equation}
where $\beta_M^{e} \approx 1.1\,\beta_M$ is the Mie \emph{extinction}
coefficient, slightly larger than the scattering coefficient because aerosols
also absorb a little light. Equation~\eqref{eq:transmittance} is evaluated per
colour channel because $\beta_R$ is.

Transmittances multiply along concatenated paths, equivalently their optical
depths add:
\begin{equation}
  T(\gamma_1 \cup \gamma_2) = T(\gamma_1)\,T(\gamma_2).
  \label{eq:multiply}
\end{equation}

\section{Phase functions}

A phase function $P(\theta)$ gives the fraction of scattered light that leaves
in a direction making angle $\theta$ with the incoming light. It is normalised
so that $\int_{4\pi} P \dd\omega = 1$. With $\mu = \cos\theta$:

\paragraph{Rayleigh.}
\begin{equation}
  P_R(\mu) = \frac{3}{16\pi}\left(1 + \mu^2\right).
  \label{eq:phase_rayleigh}
\end{equation}

\paragraph{Mie (Henyey--Greenstein).}
\begin{equation}
  P_M(\mu) = \frac{1 - g^2}{4\pi \left(1 + g^2 - 2 g \mu\right)^{3/2}},
  \qquad g \in (-1, 1).
  \label{eq:phase_hg}
\end{equation}
The asymmetry parameter $g$ controls how forward-peaked the scattering is;
atmospheric aerosols use $g \approx 0.76$.

\paragraph{Mie (Cornette--Shanks).} A more physical alternative with the same
$g$:
\begin{equation}
  P_M(\mu) = \frac{3}{8\pi}\,
             \frac{\left(1 - g^2\right)\left(1 + \mu^2\right)}
                  {\left(2 + g^2\right)\left(1 + g^2 - 2 g \mu\right)^{3/2}} .
  \label{eq:phase_cs}
\end{equation}

\section{The single-scattering integral}

Let the camera ray be $\vect{x}(t) = \vect{o} + t\,\vect{d}$ with
$\lVert\vect{d}\rVert = 1$. Let $[t_a, t_b]$ be the interval over which the ray is
inside the shell (Section~\ref{sec:shell}). Let $\vect{s}$ be the unit vector
toward the sun and $I_{\mathrm{sun}}$ the sun's irradiance at the top of the
atmosphere. The scattering angle is the same for every point on the ray:
\begin{equation}
  \mu = \vect{d} \cdot \vect{s}.
\end{equation}

Consider a point $\vect{p} = \vect{x}(t)$. Sunlight reaching $\vect{p}$ has been
attenuated by the path $\gamma_{\mathrm{sun}}(\vect{p})$ from the top of the
atmosphere to $\vect{p}$ along $-\vect{s}$. A fraction $\beta\,\rho(h)\,P(\mu)\dd t$
of it is redirected toward the camera, and that light is attenuated again by
the path $\gamma_{\mathrm{cam}}(\vect{p})$ from $\vect{p}$ back to $\vect{o}$. Summing
over the ray, for one particle type $X \in \{R, M\}$:
\begin{equation}
  L_X = I_{\mathrm{sun}}\, \beta_X\, P_X(\mu)
        \int_{t_a}^{t_b}
          \rho_X\big(h(\vect{x}(t))\big)\,
          T\big(\gamma_{\mathrm{sun}}(\vect{x}(t))\big)\,
          T\big(\gamma_{\mathrm{cam}}(\vect{x}(t))\big)
        \dd t .
  \label{eq:single_scatter}
\end{equation}
Using \eqref{eq:multiply}, the product of the two transmittances is a single
exponential of the summed optical depths, which is what the implementation
computes. The total in-scattered radiance is
\begin{equation}
  L_{\mathrm{in}} = L_R + L_M ,
  \label{eq:lin}
\end{equation}
and the transmittance used in \eqref{eq:composite} is that of the full view
path,
\begin{equation}
  T = T\big([t_a, t_b]\big) .
  \label{eq:tview}
\end{equation}

If $\gamma_{\mathrm{sun}}(\vect{p})$ intersects the planet, $\vect{p}$ is in the
planet's shadow and its contribution is zero.

\section{Ray--shell intersection}
\label{sec:shell}

For a sphere of radius $R$ centred at $\vect{c}$, substituting the ray into
$\lVert \vect{x} - \vect{c} \rVert^2 = R^2$ gives a quadratic in $t$. With
$\vect{m} = \vect{o} - \vect{c}$,
\begin{equation}
  t^2 + 2\,(\vect{m}\cdot\vect{d})\,t + \left(\vect{m}\cdot\vect{m} - R^2\right) = 0,
  \qquad
  t_{0,1} = -(\vect{m}\cdot\vect{d}) \mp
            \sqrt{(\vect{m}\cdot\vect{d})^2 - \left(\vect{m}\cdot\vect{m} - R^2\right)} .
  \label{eq:sphere}
\end{equation}
Both roots are needed, unlike a normal surface hit. The march interval is
\begin{equation}
  t_a = \max(t_0^{(\Ra)},\, 0), \qquad
  t_b = \min\!\big(t_1^{(\Ra)},\, t_{\mathrm{hit}}\big),
\end{equation}
where $t_{\mathrm{hit}}$ is the distance to the nearest surface hit (the planet
or any other object), or $\infty$ if the ray escapes. Clamping $t_a$ at zero
handles a camera that is already inside the atmosphere. The same formula with
$R = \Rp$ is used to end the light ray early when it strikes the planet.

\section{Numerical evaluation}

The integral in \eqref{eq:single_scatter} has no closed form, so it is
evaluated with a nested Riemann sum (ray marching). The view ray is divided
into $N$ segments of length $\Delta t = (t_b - t_a)/N$ with midpoints
$\vect{p}_i$; from each $\vect{p}_i$ the light ray toward the sun is divided
into $M$ segments of length $\Delta l_i$ with midpoints $\vect{q}_{ij}$.

\begin{equation}
  \tau^{\mathrm{cam}}_{X,i} = \sum_{k \le i} \rho_X(h(\vect{p}_k))\,\Delta t, \qquad
  \tau^{\mathrm{sun}}_{X,i} = \sum_{j=1}^{M} \rho_X(h(\vect{q}_{ij}))\,\Delta l_i .
\end{equation}
\begin{equation}
  A_i = \exp\!\Big(-\big[
        \beta_R\,(\tau^{\mathrm{cam}}_{R,i} + \tau^{\mathrm{sun}}_{R,i}) +
        \beta_M^{e}\,(\tau^{\mathrm{cam}}_{M,i} + \tau^{\mathrm{sun}}_{M,i})
        \big]\Big)
\end{equation}
\begin{equation}
  L_{\mathrm{in}} \approx I_{\mathrm{sun}}
  \left[
    \beta_R\, P_R(\mu) \sum_{i=1}^{N} A_i\, \rho_R(h(\vect{p}_i))\,\Delta t
  + \beta_M\, P_M(\mu) \sum_{i=1}^{N} A_i\, \rho_M(h(\vect{p}_i))\,\Delta t
  \right]
\end{equation}
\begin{equation}
  T \approx \exp\!\Big(-\big[\beta_R\,\tau^{\mathrm{cam}}_{R,N} +
                            \beta_M^{e}\,\tau^{\mathrm{cam}}_{M,N}\big]\Big)
\end{equation}

Note that the Mie \emph{extinction} coefficient $\beta_M^e$ appears inside the
attenuation while the Mie \emph{scattering} coefficient $\beta_M$ multiplies the
in-scatter sum.

\begin{algorithm}[h]
\caption{Single-scattering march for one camera ray}
\begin{algorithmic}[1]
\Require ray $(\vect{o}, \vect{d})$, interval $[t_a, t_b]$, sun direction $\vect{s}$
\State $\tau^{\mathrm{cam}}_R, \tau^{\mathrm{cam}}_M, S_R, S_M \gets 0$
\State $\Delta t \gets (t_b - t_a)/N$;\quad $\mu \gets \vect{d}\cdot\vect{s}$
\For{$i = 1 \ldots N$}
  \State $\vect{p} \gets \vect{o} + \big(t_a + (i - \tfrac12)\Delta t\big)\vect{d}$;\quad $h \gets \lVert\vect{p} - \vect{c}\rVert - \Rp$
  \State $d_R \gets \rho_R(h)\,\Delta t$;\quad $d_M \gets \rho_M(h)\,\Delta t$
  \State $\tau^{\mathrm{cam}}_R \mathrel{+}= d_R$;\quad $\tau^{\mathrm{cam}}_M \mathrel{+}= d_M$
  \State find $t_{\mathrm{exit}}$ where the ray $(\vect{p}, \vect{s})$ leaves the shell of radius $\Ra$
  \If{the ray $(\vect{p}, \vect{s})$ hits the sphere of radius $\Rp$} \textbf{continue} \Comment{in planet shadow} \EndIf
  \State $\tau^{\mathrm{sun}}_R, \tau^{\mathrm{sun}}_M \gets 0$;\quad $\Delta l \gets t_{\mathrm{exit}}/M$
  \For{$j = 1 \ldots M$}
    \State $\vect{q} \gets \vect{p} + (j - \tfrac12)\Delta l\,\vect{s}$;\quad $h' \gets \lVert\vect{q} - \vect{c}\rVert - \Rp$
    \State $\tau^{\mathrm{sun}}_R \mathrel{+}= \rho_R(h')\,\Delta l$;\quad $\tau^{\mathrm{sun}}_M \mathrel{+}= \rho_M(h')\,\Delta l$
  \EndFor
  \State $A \gets \exp\!\big(-[\beta_R(\tau^{\mathrm{cam}}_R + \tau^{\mathrm{sun}}_R) + \beta_M^e(\tau^{\mathrm{cam}}_M + \tau^{\mathrm{sun}}_M)]\big)$
  \State $S_R \mathrel{+}= A\,d_R$;\quad $S_M \mathrel{+}= A\,d_M$
\EndFor
\State $L_{\mathrm{in}} \gets I_{\mathrm{sun}}\big(\beta_R P_R(\mu) S_R + \beta_M P_M(\mu) S_M\big)$
\State $T \gets \exp\!\big(-[\beta_R \tau^{\mathrm{cam}}_R + \beta_M^e \tau^{\mathrm{cam}}_M]\big)$
\State \Return $(L_{\mathrm{in}}, T)$
\end{algorithmic}
\end{algorithm}

\section{Tone mapping}

$L_{\mathrm{in}}$ is high dynamic range and routinely exceeds $1$. Before
quantising to 8 bits each channel is passed through an exposure curve such as
\begin{equation}
  L' = 1 - \exp(-E\,L),
  \label{eq:tonemap}
\end{equation}
with exposure $E$ chosen per scene, followed by an optional gamma
$L'' = (L')^{1/2.2}$.

\section{Reference constants}

Values for an Earth-like atmosphere, expressed in kilometres so that
single-precision floats retain adequate precision at planetary scale.

\begin{center}
\begin{tabular}{llr}
\toprule
Symbol & Meaning & Value \\
\midrule
$\Rp$ & planet radius & \SI{6371}{\kilo\metre} \\
$\Ra$ & atmosphere radius & \SI{6471}{\kilo\metre} \\
$H_R$ & Rayleigh scale height & \SI{8.0}{\kilo\metre} \\
$H_M$ & Mie scale height & \SI{1.2}{\kilo\metre} \\
$\beta_R$ & Rayleigh scattering (R, G, B) & $(5.8,\ 13.5,\ 33.1)\times 10^{-3}\ \si{\per\kilo\metre}$ \\
$\beta_M$ & Mie scattering & $21\times 10^{-3}\ \si{\per\kilo\metre}$ \\
$\beta_M^e$ & Mie extinction & $1.1\,\beta_M$ \\
$g$ & Mie asymmetry & $0.76$ \\
$I_{\mathrm{sun}}$ & sun irradiance (before exposure) & $\approx 20$ \\
$N$ & view-ray samples & 16 (start), 64 (final) \\
$M$ & light-ray samples & 8 (start), 16 (final) \\
\bottomrule
\end{tabular}
\end{center}

\section{Qualitative predictions}

The formulas above account for the familiar appearance of the sky:
\begin{itemize}
  \item \textbf{Blue daytime sky.} $\beta_R$ is largest in blue, so
        \eqref{eq:single_scatter} redirects more blue light toward the viewer,
        and the short overhead path leaves the view transmittance near $1$.
  \item \textbf{Red sunsets and terminator.} Light reaching air near the viewer
        has crossed a long, low, dense path along $-\vect{s}$, so
        $T(\gamma_{\mathrm{sun}})$ in \eqref{eq:transmittance} is strongly
        depleted in blue and green.
  \item \textbf{Bright limb from orbit.} A view ray tangent to the surface has a
        large $\tau^{\mathrm{cam}}$ and therefore a long column of scatterers,
        increasing $L_{\mathrm{in}}$.
  \item \textbf{Solar halo and horizon haze.} $P_M(\mu)$ in \eqref{eq:phase_hg}
        grows sharply as $\mu \to 1$, and $\rho_M$ is concentrated near the
        ground because $H_M < H_R$.
\end{itemize}

The single-scattering assumption under-estimates brightness on the night side
of the terminator and in twilight, where most visible light has scattered more
than once.

\section*{References}
\begin{itemize}
  \item T.~Nishita, T.~Sirai, K.~Tadamura, E.~Nakamae.
        \emph{Display of the Earth Taking into Account Atmospheric Scattering}.
        SIGGRAPH 1993.
  \item Scratchapixel, \emph{Simulating the Colors of the Sky}.
        \url{https://www.scratchapixel.com/lessons/procedural-generation-virtual-worlds/simulating-sky}
  \item S.~O'Neil. \emph{Accurate Atmospheric Scattering}. GPU Gems 2, Chapter 16.
\end{itemize}

\end{document}
